\section{おわりに}
本研究では，環境変化に対して迅速に適応する強化学習の実現に向け，
方策再利用（Policy Reuse）の枠組みにおける方策選択に着目し，
既存研究で用いられているボルツマン選択に加えて 
$\varepsilon$-greedy および UCB を導入した比較実験を行った．
さらに，報酬設計を2種類（ゴール到達報酬のみ，
およびステップ負報酬＋壁衝突ペナルティ）用意し，
報酬構造の変更が各方策選択手法の挙動と性能差に与える影響も検証した．
その結果，いずれの報酬設計においても
方策再利用を用いるPRQ系列はBaseline（通常のQ学習）を
一貫して上回り，転移（再利用）による学習効率向上が確認できた．
一方で，PRQ内の方策選択手法の優劣はタスクや報酬設計に依存して
変化し，単一の手法が常に最良となるわけではないことが示唆された．

今後の課題として，
より多様な迷路構造や報酬スケールに対する再現性検証に加え，
グリッドサーチに依存しない形で探索パラメータを自動調整する手法の
導入が挙げられる．
また，本研究は表形式のQ学習を前提としたが，
より大規模な状態空間や実問題への適用を想定すると，
関数近似や深層強化学習への拡張が重要となる．
加えて，複数のソース環境から得られた方策を
ライブラリに蓄積する状況では，
方策の選択・管理戦略そのものが性能に直結するため，
方策集合の構成方法や選択基準の高度化も検討すべきである．
これらの検討を通じて，環境変化を伴う実運用場面においても
頑健かつ効率的に学習できる転移強化学習手法の確立を目指す．